% Abhakseite „Das kann ich“ – Zone nur im Gesamt (\mitzone)
%% TODO Ich-kann-Sätze: Platzhalter wie in den Aufgabendateien
\begin{abhakseite}
\ifdefined\mitzone
\abhakgruppe{Kennst du schon}
\abhak{1}{Punkte im räumlichen Koordinatensystem lesen und Lagen aus Koordinaten deuten}
\abhak{2}{Verbindungsvektor bilden, Vielfache erkennen, den Betrag berechnen und Vektoren normieren}
\abhak{3}{Skalarprodukt berechnen und Orthogonalität erkennen}
\abhak{4}{Lineare Gleichungen lösen und Widersprüche erkennen}
\abhak{5}{Kollinearität prüfen}
\abhak{6}{Die Gerade als Gleichung in der Ebene lesen (y = m · x + n, Anstieg)}
\abhak{7}{Prozentrechnung und Maßstab}
\abhak{8}{Lineare Gleichungen lösen und Widersprüche erkennen – Fehler finden}
\abhak{9}{Lineare Gleichungen lösen und Widersprüche erkennen – selbst rechnen}
\fi
\abhakgruppe{Einheit 1 · Geradengleichung aufstellen und lesen}
\abhak{10}{„Gerade oder Strecke?“}
\abhak{11}{Geradengleichung aufstellen und lesen}
\abhak{12}{Geradengleichung aufstellen und lesen – Fehler finden, Begründen, Anwendung, Darstellungswechsel}
\abhakgruppe{Einheit 2 · Punktprobe und Punkte auf der Geraden}
\abhak{13}{Punktprobe und Punkte auf der Geraden}
\abhak{14}{Punkt auf einer Geraden mit vorgegebenem Abstand zu einem festen Punkt über eine quadratische Gleichung bestimmen}
\abhak{15}{Punktprobe und Punkte auf der Geraden – Fehler finden, Begründen, Anwendung}
\abhakgruppe{Einheit 3 · Lage als Anhang}
\abhak{16}{Lage als Anhang}
\abhak{17}{Lage als Anhang – Fehler finden, Begründen}
\abhakgruppe{Einheit 4 · Sachgeraden}
\abhak{18}{Sachgeraden}
\abhak{19}{Horizontalabstand aus vorgegebener Neigung in Prozent und Höhendifferenz berechnen}
\abhak{20}{Sachgeraden – Fehler finden, Begründen, Anwendung}
\end{abhakseite}
